%% kapa: exam=1 round=20 year=2009 period=2 form=A subject=appraisal_law subject_ko=부동산관계법규 seq=14 number=54
%% source: Q-Net 감정평가사 기출문제 (cst00309) artlSeq=5024178 file=extracted/1008/5024178/2009년 제20회 감정평가사 1차시험 기출문제/2교시.hwp
%% conversion: kordoc 4.24.51 -> tools/gen1.py (auto); review=proofread-sonnet; flags=none
\begin{question}{54}
국토의 계획 및 이용에 관한 법령상 비용부담에 관한 설명으로 옳지 \uline{않은} 것은?
\begin{choices}
  \item 도시계획시설사업에 관한 비용에 관하여 이 법 또는 다른 법률에 특별한 규정이 없는 한 행정청이 아닌 자가 시행하는 경우에는 그 자가 부담함을 원칙으로 한다.
  \item 국토해양부장관 또는 시ㆍ도지사는 그가 시행한 도시계획시설사업으로 인하여 현저히 이익을 받는 시ㆍ도가 있는 때에는 당해 도시계획시설사업에 소요된 비용의 전부를 그 이익을 받는 시ㆍ도에 부담시킬 수 있다.
  \item 도시관리계획의 입안을 제안 받은 자는 제안자와 협의하여 제안된 도시관리계획의 입안 및 결정에 필요한 비용의 전부 또는 일부를 제안자에게 부담시킬 수 있다.
  \item 행정청인 도시계획시설사업의 시행자는 공동구를 설치하는 경우 다른 법률에 따라 그 공동구에 수용되어야 할 시설을 설치할 의무가 있는 자에게 공동구의 설치에 필요한 비용을 부담시킬 수 있다.
  \item 행정청이 시행하는 도시계획시설사업에 드는 비용은 대통령령으로 정하는 바에 따라 그 비용의 전부 또는 일부를 국가예산에서 보조하거나 융자할 수 있다.
\end{choices}
\end{question}
\answer{2}
